% Generated draft: requires mathematical and source review.
\documentclass[11pt]{article}
\usepackage[margin=25mm]{geometry}
\usepackage{fontspec}
\setmainfont{Latin Modern Roman}
\usepackage{amsmath,amssymb,mathtools,unicode-math}
\setmathfont{Latin Modern Math}
\usepackage{enumitem,needspace,xurl,hyperref}
\setlength{\parindent}{0pt}
\setlength{\parskip}{5pt}
\setlength{\emergencystretch}{3em}
\newcommand{\N}{\mathbb{N}}
\newcommand{\Z}{\mathbb{Z}}
\newcommand{\Q}{\mathbb{Q}}
\newcommand{\R}{\mathbb{R}}
\newcommand{\C}{\mathbb{C}}
\newcommand{\F}{\mathbb{F}}
\newcommand{\A}{\mathbb{A}}
\newcommand{\PP}{\mathbb{P}}
\newcommand{\E}{\mathbb{E}}
\newcommand{\eps}{\varepsilon}
\newcommand{\dd}{\,\mathrm{d}}
\begin{document}
\section*{UnsolvedMath ID: 30003593}
\textbf{Source problem code:} catalogue entry
\subsection{Definitions and mathematical statement}
In the usual projective setting, let \(Z_1(X)_{\R}\) denote the real vector space of finite real linear combinations of integral curves on \(X\). Two such cycles \(\alpha\) and \(\beta\) are numerically equivalent if
\[
D\cdot\alpha=D\cdot\beta\qquad\text{for every Cartier divisor }D\text{ on }X.
\]
The space of numerical curve classes is
\[
N_1(X)=Z_1(X)_{\R}/\!\equiv_{\mathrm{num}}.
\]
A convex body means a nonempty compact convex subset of a finite-dimensional real vector space. These definitions explain the customary notation; the source does not specify the class of varieties \(X\), a ground field, a dimension, or regularity assumptions.

\textbf{Source question.} Can the positivity theory for curve classes in \(N_1(X)\) be expressed through Newton--Okounkov-type convex bodies?

\textbf{Unresolved formulation.} The supplied question asks for a theory rather than a specified mathematical assertion. It does not identify the curve classes to which bodies should be assigned, the positivity properties or numerical invariants to be recovered, the auxiliary choices allowed, or the properties that make a construction ``Newton--Okounkov-type.'' Consequently, no precise universally quantified existence conjecture can be extracted without additional hypotheses and a definition of the required correspondence. The projective setting used above to explain \(N_1(X)\) is contextual and is not a hypothesis explicitly supplied by the record.

The accompanying literature assessment concerns convex-body models for higher asymptotic cohomology functions, whereas the statement concerns positivity of numerical curve classes. The record supplies no equivalence or other connection that identifies these two targets; they must remain distinct pending review.
\subsection{Short English statement}
The question asks whether Newton--Okounkov-type convex bodies can describe positivity of numerical curve classes. The permitted varieties, relevant positivity notions, and required properties of the construction are unspecified. The accompanying discussion of higher asymptotic cohomology does not determine the intended curve-class statement.
\subsection{Sources}
\begin{itemize}
\item Local UnsolvedMath registry, permanent ID \(30003593\), title \emph{Newton--Okounkov Geometry for Positive Curve Classes}, exact source problem code \texttt{catalogue entry}; collection: \emph{Oberwolfach Reports Open Problems}. No individual proposer is supplied. The proposed year and year tag are 2017. The supplied status is \texttt{open}; the supplied statement status is \texttt{corrected\_verified}, with an explanation that notation, typography, and sentence structure were repaired while preserving mathematical content. These are registry labels, not independent verification.
\item Supplied report citation: \emph{Mini-Workshop: Positivity in Higher-dimensional Geometry: Higher-codimensional Cycles and Newton--Okounkov Bodies} (2018), \emph{Oberwolfach Reports}, pp.\ 2631--2657, DOI 10.4171/owr/2017/43. Supplied source URL: \url{https://doi.org/10.4171/owr/2017/43}. The literature-source list also gives the original OWR report as \url{https://doi.org/10.4171/OWR/2017/43}. The supplied background describes these reports as records of mathematical workshops and their open-problem discussions.
\item Supplied literature reference: K\"uronya, asymptotic cohomological functions, \url{https://arxiv.org/abs/math/0501491}.
\item Supplied literature reference: Hering--K\"uronya--Payne, toric asymptotic cohomology, \url{https://arxiv.org/abs/math/0501104}.
\item Supplied literature assessment, checked on exactly \texttt{2026-08-21}: status \texttt{open}. The assessment states that Newton--Okounkov bodies give convex interpretations of volume and several restricted invariants, that higher asymptotic cohomology has only special-case formulas such as toric and surface results, and that no general Newton--Okounkov-type convex-body model for higher cohomology functions is known. This is supplied dated curation derived from the OWR material, not independently verified. Its time-sensitive status and cited claims require independent checking, and its relationship to the curve-class question is unresolved.
\end{itemize}
\end{document}
